% GENERATED FILE. Edit canonical lessons, then run npm run generate:books.
% canonical-source-hash: fnv1a64:98dc4f3033fa81d8
% canonical-lessons: KA-C219-samanya, KA-C219-mondu, KA-C219-nunupu, KA-C219-oratu, KA-C219-jigutu

\chapter{Ordinary, Blunt, Smooth, Rough, Sticky}
\label{ch:ka-ordinary-blunt-smooth-rough-sticky}
\hlchaptermodality{219}

\begin{chapteropening}
\textbf{By the end of this chapter:} \emph{I can say ordinary, blunt, smooth, rough and sticky.}

Complete the last lesson of chapter 219: I can say ordinary, blunt, smooth, rough and sticky.
\end{chapteropening}

\section[\texorpdfstring{\kn{ಸಾಮಾನ್ಯ} (sāmānya)}{ಸಾಮಾನ್ಯ (sāmānya)}]{\kn{ಸಾಮಾನ್ಯ} (sāmānya) — ordinary, common}

\label{lesson:KA-C219-samanya}



\begin{quote}
Before the new one: say the Kannada for honest, then the Kannada for generous.
\end{quote}

\subsection*{You'll want to know: \kn{ಸಾಮಾನ್ಯ}}
\textbf{\kn{ಸಾಮಾನ್ಯ}} — \emph{sāmānya} — \textquotedblleft{}ordinary, common\textquotedblright{}.

\begin{grammarlens}[title={What you've built}]
One more word for this chapter.
\end{grammarlens}

\subsection*{Your turn}
\begin{itemize}
\raggedright
  \item \emph{Say it:} \emph{sāmānya}
  \item \emph{Say it:} \emph{sāmānya}, once more
  \item \emph{Say it:} say \emph{sāmānya}
  \item \emph{Recall:} say the Kannada for honest, then the Kannada for generous, then say \emph{sāmānya} again
\end{itemize}

\begin{tcolorbox}[breakable,colback=teal!4,colframe=teal!35!black,title={Before you move on}]
What does \kn{ಸಾಮಾನ್ಯ} mean? (\textquotedblleft{}Ordinary, common\textquotedblright{}.) Say it once more.
\end{tcolorbox}
\section[\texorpdfstring{\kn{ಮೊಂಡು} (moṇḍu)}{ಮೊಂಡು (moṇḍu)}]{\kn{ಮೊಂಡು} (moṇḍu) — blunt}

\label{lesson:KA-C219-mondu}



\begin{quote}
Before the new one: say the Kannada for generous, then the Kannada for ordinary.
\end{quote}

\subsection*{You'll want to know: \kn{ಮೊಂಡು}}
\textbf{\kn{ಮೊಂಡು}} — \emph{moṇḍu} — \textquotedblleft{}blunt\textquotedblright{}.

\begin{grammarlens}[title={What you've built}]
One more word for this chapter.
\end{grammarlens}

\subsection*{Your turn}
\begin{itemize}
\raggedright
  \item \emph{Say it:} \emph{moṇḍu}
  \item \emph{Say it:} \emph{moṇḍu}, once more
  \item \emph{Say it:} say \emph{moṇḍu}
  \item \emph{Recall:} say the Kannada for generous, then the Kannada for ordinary, then say \emph{moṇḍu} again
\end{itemize}

\begin{tcolorbox}[breakable,colback=teal!4,colframe=teal!35!black,title={Before you move on}]
What does \kn{ಮೊಂಡು} mean? (\textquotedblleft{}Blunt\textquotedblright{}.) Say it once more.
\end{tcolorbox}
\section[\texorpdfstring{\kn{ನುಣುಪು} (nuṇupu)}{ನುಣುಪು (nuṇupu)}]{\kn{ನುಣುಪು} (nuṇupu) — smooth}

\label{lesson:KA-C219-nunupu}



\begin{quote}
Before the new one: say the Kannada for ordinary, then the Kannada for blunt.
\end{quote}

\subsection*{You'll want to know: \kn{ನುಣುಪು}}
\textbf{\kn{ನುಣುಪು}} — \emph{nuṇupu} — \textquotedblleft{}smooth\textquotedblright{}.

\begin{grammarlens}[title={What you've built}]
One more word for this chapter.
\end{grammarlens}

\subsection*{Your turn}
\begin{itemize}
\raggedright
  \item \emph{Say it:} \emph{nuṇupu}
  \item \emph{Say it:} \emph{nuṇupu}, once more
  \item \emph{Say it:} say \emph{nuṇupu}
  \item \emph{Recall:} say the Kannada for ordinary, then the Kannada for blunt, then say \emph{nuṇupu} again
\end{itemize}

\begin{tcolorbox}[breakable,colback=teal!4,colframe=teal!35!black,title={Before you move on}]
What does \kn{ನುಣುಪು} mean? (\textquotedblleft{}Smooth\textquotedblright{}.) Say it once more.
\end{tcolorbox}
\section[\texorpdfstring{\kn{ಒರಟು} (oraṭu)}{ಒರಟು (oraṭu)}]{\kn{ಒರಟು} (oraṭu) — rough, coarse}

\label{lesson:KA-C219-oratu}



\begin{quote}
Before the new one: say the Kannada for blunt, then the Kannada for smooth.
\end{quote}

\subsection*{You'll want to know: \kn{ಒರಟು}}
\textbf{\kn{ಒರಟು}} — \emph{oraṭu} — \textquotedblleft{}rough, coarse\textquotedblright{}.

\begin{grammarlens}[title={What you've built}]
One more word for this chapter.
\end{grammarlens}

\subsection*{Your turn}
\begin{itemize}
\raggedright
  \item \emph{Say it:} \emph{oraṭu}
  \item \emph{Say it:} \emph{oraṭu}, once more
  \item \emph{Say it:} say \emph{oraṭu}
  \item \emph{Recall:} say the Kannada for blunt, then the Kannada for smooth, then say \emph{oraṭu} again
\end{itemize}

\begin{tcolorbox}[breakable,colback=teal!4,colframe=teal!35!black,title={Before you move on}]
What does \kn{ಒರಟು} mean? (\textquotedblleft{}Rough, coarse\textquotedblright{}.) Say it once more.
\end{tcolorbox}
\section[\texorpdfstring{\kn{ಜಿಗುಟು} (jiguṭu)}{ಜಿಗುಟು (jiguṭu)}]{\kn{ಜಿಗುಟು} (jiguṭu) — sticky}

\label{lesson:KA-C219-jigutu}



\begin{quote}
Before the new one: say the Kannada for smooth, then the Kannada for rough.
\end{quote}

\subsection*{You'll want to know: \kn{ಜಿಗುಟು}}
\textbf{\kn{ಜಿಗುಟು}} — \emph{jiguṭu} — \textquotedblleft{}sticky\textquotedblright{}.

\begin{grammarlens}[title={What you've built}]
One more word for this chapter.
\end{grammarlens}

\subsection*{Your turn}
\begin{itemize}
\raggedright
  \item \emph{Say it:} \emph{jiguṭu}
  \item \emph{Say it:} \emph{jiguṭu}, once more
  \item \emph{Say it:} say \emph{jiguṭu}
  \item \emph{Recall:} say the Kannada for smooth, then the Kannada for rough, then say \emph{jiguṭu} again
\end{itemize}

\begin{tcolorbox}[breakable,colback=teal!4,colframe=teal!35!black,title={Before you move on}]
What does \kn{ಜಿಗುಟು} mean? (\textquotedblleft{}Sticky\textquotedblright{}.) Say it once more.
\end{tcolorbox}
